\documentclass[10pt,a4paper]{amsart}
\usepackage{amsmath,amssymb,amsthm,graphicx}
\usepackage{verbatim}
\newenvironment{reference}{\comment}{\endcomment}
%\newenvironment{reference}{}{}
\newenvironment{slogan}{\comment}{\endcomment}
\newenvironment{history}{\comment}{\endcomment}
% For commutative diagrams we use Xy-pic
\usepackage[all]{xy}

% We use 2cell for 2-commutative diagrams.
\xyoption{2cell}
% Theorem environments.
%
\theoremstyle{plain}
\newtheorem{theorem}[subsection]{Theorem}
\newtheorem{proposition}[subsection]{Proposition}
\newtheorem{lemma}[subsection]{Lemma}

\theoremstyle{definition}
\newtheorem{definition}[subsection]{Definition}
\newtheorem{example}[subsection]{Example}
\newtheorem{exercise}[subsection]{Exercise}
\newtheorem{situation}[subsection]{Situation}

\theoremstyle{remark}
\newtheorem{remark}[subsection]{Remark}
\newtheorem{remarks}[subsection]{Remarks}

\numberwithin{equation}{subsection}

% Macros
%
\def\lim{\mathop{\mathrm{lim}}\nolimits}
\def\colim{\mathop{\mathrm{colim}}\nolimits}
\def\Spec{\mathop{\mathrm{Spec}}}
\def\Hom{\mathop{\mathrm{Hom}}\nolimits}
\def\Ext{\mathop{\mathrm{Ext}}\nolimits}
\def\SheafHom{\mathop{\mathcal{H}\!\mathit{om}}\nolimits}
\def\SheafExt{\mathop{\mathcal{E}\!\mathit{xt}}\nolimits}
\def\Sch{\mathit{Sch}}
\def\Mor{\mathop{\mathrm{Mor}}\nolimits}
\def\Ob{\mathop{\mathrm{Ob}}\nolimits}
\def\Sh{\mathop{\mathit{Sh}}\nolimits}
\def\NL{\mathop{N\!L}\nolimits}
\def\CH{\mathop{\mathrm{CH}}\nolimits}
\def\proetale{{pro\text{-}\acute{e}tale}}
\def\etale{{\acute{e}tale}}
\def\QCoh{\mathit{QCoh}}
\def\Ker{\mathop{\mathrm{Ker}}}
\def\Im{\mathop{\mathrm{Im}}}
\def\Coker{\mathop{\mathrm{Coker}}}
\def\Coim{\mathop{\mathrm{Coim}}}

% Boxtimes
%
\DeclareMathSymbol{\boxtimes}{\mathbin}{AMSa}{"02}

%
% Macros for moduli stacks/spaces
%
\def\QCohstack{\mathcal{QC}\!\mathit{oh}}
\def\Cohstack{\mathcal{C}\!\mathit{oh}}
\def\Spacesstack{\mathcal{S}\!\mathit{paces}}
\def\Quotfunctor{\mathrm{Quot}}
\def\Hilbfunctor{\mathrm{Hilb}}
\def\Curvesstack{\mathcal{C}\!\mathit{urves}}
\def\Polarizedstack{\mathcal{P}\!\mathit{olarized}}
\def\Complexesstack{\mathcal{C}\!\mathit{omplexes}}
% \Pic is the operator that assigns to X its picard group, usage \Pic(X)
% \Picardstack_{X/B} denotes the Picard stack of X over B
% \Picardfunctor_{X/B} denotes the Picard functor of X over B
\def\Pic{\mathop{\mathrm{Pic}}\nolimits}
\def\Picardstack{\mathcal{P}\!\mathit{ic}}
\def\Picardfunctor{\mathrm{Pic}}
\def\Deformationcategory{\mathcal{D}\!\mathit{ef}}
\usepackage{cleveref}
\crefname{originaldefinition}{original definition}{original definitions}
\Crefname{originaldefinition}{Original definition}{Original definitions}
\makeatletter
\newlabel{divisors-definition-relative-effective-Cartier-divisor}{{062T}{source}{}{sourcepointer-1}{}}
\newlabel{divisors-definition-relative-effective-Cartier-divisor@cref}{{[originaldefinition][0][]062T}{[1][0][]source}{}{}{}}
\makeatother
\crefname{originallemma}{original lemma}{original lemmas}
\Crefname{originallemma}{Original lemma}{Original lemmas}
\makeatletter
\newlabel{divisors-lemma-relative-Cartier}{{056Q}{source}{}{sourcepointer-2}{}}
\newlabel{divisors-lemma-relative-Cartier@cref}{{[originallemma][0][]056Q}{[1][0][]source}{}{}{}}
\makeatother
\crefname{originallemma}{original lemma}{original lemmas}
\Crefname{originallemma}{Original lemma}{Original lemmas}
\makeatletter
\newlabel{lemma-easy-gysin}{{02TB}{source}{}{sourcepointer-3}{}}
\newlabel{lemma-easy-gysin@cref}{{[originallemma][0][]02TB}{[1][0][]source}{}{}{}}
\makeatother
\crefname{originallemma}{original lemma}{original lemmas}
\Crefname{originallemma}{Original lemma}{Original lemmas}
\makeatletter
\newlabel{lemma-support-cap-effective-Cartier}{{02T9}{source}{}{sourcepointer-4}{}}
\newlabel{lemma-support-cap-effective-Cartier@cref}{{[originallemma][0][]02T9}{[1][0][]source}{}{}{}}
\makeatother
\crefname{originalreference}{original reference}{original references}
\Crefname{originalreference}{Original reference}{Original references}
\makeatletter
\newlabel{situation-setup}{{02QL}{source}{}{sourcepointer-5}{}}
\newlabel{situation-setup@cref}{{[originalreference][0][]02QL}{[1][0][]source}{}{}{}}
\makeatother
\title{Gysin and Chern class expressions for a relative effective Cartier divisor}
\author{The Stacks Project Authors}
\date{stacks-project:chow.tex:local-snapshot-20261012; GNU FDL 1.2 or later}
\begin{document}
\maketitle
\noindent\textit{Editorial note.} All delivered mathematical text is the author's own, copied from the pinned Stacks Project chapter file chow.tex; no statement or proof step was rewritten, reordered or removed.\par\medskip
\noindent\textit{Editorial note.} Preamble provenance: the pinned chapter file opens with \textbackslash{}input\{preamble\} and carries no preamble text of its own, so the project's single Stacks-wide preamble staged at sources/stacks-chow/preamble.tex is excerpted verbatim as three spans, lines 8-12 (the verbatim package together with the project's own reference, slogan and history display-environment definitions, including the author's commented-out older form of the reference environment), lines 14-18 (Xy-pic setup) and lines 154-217 (the native theorem declarations and mathematical macros). The lines 8-12 span is required because a counted statement or its surrounding text may use those display environments, which the project typesets as comments. Nothing was added, removed or redefined; the wrapper class and the amsmath/amssymb/amsthm/graphicx base packages are supplied by the assembler.\par\medskip
\noindent\textit{Editorial note.} Same-chapter pointers are located at their \textbackslash{}label lines inside the pinned chapter file; cross-chapter pointers are located in the project's own tag list staged in the snapshot (sources/stacks-chow/tags.txt), at the verbatim lines carrying their canonical tags.\par\medskip
\noindent\textit{Allowed standard prerequisites (editorial).}\begin{itemize}
\item A relative effective Cartier divisor restricts to an effective Cartier divisor on each fibre product, and the easy Gysin computation. Divisors, label divisors-lemma-relative-Cartier (proof line 5619); Same chapter chow.tex, label lemma-easy-gysin (proof line 5621, claimed as S14-algebra-easy-gysin). The first shows that the intersection of the divisor with the preimage of a closed subscheme is again an effective Cartier divisor, and the second computes the Gysin pullback of the resulting cycle class.
\item The Chern class expression for an effective Cartier divisor. Same chapter chow.tex, label lemma-support-cap-effective-Cartier, invoked at proof line 5634. Deduces the second formula from the first by expressing the cap with the first Chern class through the pushforward along the divisor.
\end{itemize}\medskip
\ifdefined\hypertarget\hypertarget{sourcepointer-1}{}\fi
\noindent\textit{Source reference.} Original definition 062T: sources/stacks-chow/tags.txt, source lines 7807--7807. Navigation pointer for \textbackslash{}ref\{divisors-definition-relative-effective-Cartier-divisor\} used inside the counted span (line 5604); the argument that consumes it is retained in the span.\par\medskip
\ifdefined\hypertarget\hypertarget{sourcepointer-2}{}\fi
\noindent\textit{Source reference.} Original lemma 056Q: sources/stacks-chow/tags.txt, source lines 6718--6718. Navigation pointer for \textbackslash{}ref\{divisors-lemma-relative-Cartier\} used inside the counted span (line 5619); the argument that consumes it is retained in the span.\par\medskip
\ifdefined\hypertarget\hypertarget{sourcepointer-3}{}\fi
\noindent\textit{Source reference.} Original lemma 02TB: sources/stacks-chow/pinned-original-source.tex, source lines 5170--5170. Navigation pointer for \textbackslash{}ref\{lemma-easy-gysin\} used inside the counted span (line 5621); the argument that consumes it is retained in the span.\par\medskip
\ifdefined\hypertarget\hypertarget{sourcepointer-4}{}\fi
\noindent\textit{Source reference.} Original lemma 02T9: sources/stacks-chow/pinned-original-source.tex, source lines 5138--5138. Navigation pointer for \textbackslash{}ref\{lemma-support-cap-effective-Cartier\} used inside the counted span (line 5634); the argument that consumes it is retained in the span.\par\medskip
\ifdefined\hypertarget\hypertarget{sourcepointer-5}{}\fi
\noindent\textit{Source reference.} Original reference 02QL: sources/stacks-chow/pinned-original-source.tex, source lines 1372--1372. Navigation pointer for \textbackslash{}ref\{situation-setup\} used inside the counted span (line 5599); the argument that consumes it is retained in the span.\par\medskip
\begin{lemma}
\label{lemma-relative-effective-cartier}
Let $(S, \delta)$ be as in Situation \ref{situation-setup}.
Let $X$, $Y$ be locally of finite type over $S$.
Let $p : X \to Y$ be a flat morphism of relative dimension $r$.
Let $i : D \to X$ be a relative effective Cartier divisor
(Divisors, Definition
\ref{divisors-definition-relative-effective-Cartier-divisor}).
Let $\mathcal{L} = \mathcal{O}_X(D)$.
For any $\alpha \in \CH_{k + 1}(Y)$ we have
$$
i^*p^*\alpha = (p|_D)^*\alpha
$$
in $\CH_{k + r}(D)$ and
$$
c_1(\mathcal{L}) \cap p^*\alpha = i_* ((p|_D)^*\alpha)
$$
in $\CH_{k + r}(X)$.
\end{lemma}

\begin{proof}
Let $W \subset Y$ be an integral closed subscheme of $\delta$-dimension
$k + 1$. By Divisors, Lemma \ref{divisors-lemma-relative-Cartier}
we see that $D \cap p^{-1}W$ is an effective
Cartier divisor on $p^{-1}W$. By Lemma \ref{lemma-easy-gysin}
we get the first equality in
$$
i^*[p^{-1}W]_{k + r + 1} =
[D \cap p^{-1}W]_{k + r} =
[(p|_D)^{-1}(W)]_{k + r}.
$$
and the second because $D \cap p^{-1}(W) = (p|_D)^{-1}(W)$ as schemes.
Since by definition $p^*[W] = [p^{-1}W]_{k + r + 1}$ we see that
$i^*p^*[W] = (p|_D)^*[W]$ as cycles. If $\alpha = \sum m_j[W_j]$ is a
general $k + 1$ cycle, then we get
$i^*\alpha = \sum m_j i^*p^*[W_j] = \sum m_j(p|_D)^*[W_j]$ as cycles.
This proves then first equality. To deduce the second from the
first apply Lemma \ref{lemma-support-cap-effective-Cartier}.
\end{proof}

\end{document}
